\thispagestyle{empty}

% TÍTULO DO RESUMO
\begin{center}
    \textbf{\large RESUMO}
\end{center}

\vspace{14pt}

% CORPO DO RESUMO EM PARÁGRAFO ÚNICO E ESPAÇAMENTO SIMPLES (ABNT NBR 6028)
\begin{singlespace}
\noindent\justifying
O presente trabalho apresenta o desenvolvimento do WeFIND, uma aplicação móvel colaborativa projetada para conectar tutores e a comunidade na localização, no resgate e na divulgação de animais de estimação perdidos e encontrados, além de incentivar a prática da adoção responsável. A concepção do projeto baseou-se em uma pesquisa prática com tutores e na análise comparativa de plataformas similares de mercado, identificando dificuldades enfrentadas nas buscas, como a dependência de cartazes impressos e grupos de redes sociais, e estruturando os requisitos do software por meio de diagramas de modelagem UML. A proposta integra recursos de geolocalização para visualização de publicações em mapa interativo, busca por características físicas do animal, canal de mensagens privativo para contato sem exposição de números de telefone e geração de cartazes padronizados contendo código QR para compartilhamento ou impressão. O desenvolvimento tecnológico utiliza o ecossistema React Native e Expo no ambiente móvel e a plataforma Supabase com banco de dados relacional PostgreSQL para os serviços de autenticação, armazenamento e sincronização de dados. Como contribuição, o trabalho apresenta uma solução centralizada para apoiar a colaboração comunitária na busca por animais, cuja implementação será submetida a testes funcionais, de usabilidade e de desempenho.

\vspace{18pt}

\noindent
\textbf{Palavras-chave:} Animais perdidos. Aplicativos móveis. Geolocalização. React Native. Supabase. Sistemas colaborativos.
\end{singlespace}

\clearpage
